\documentclass[a4paper, 12pt]{article}

\usepackage{verse, gmverse}
\usepackage[utf8]{inputenc}
\usepackage[T1]{fontenc}
\usepackage{textcomp}
\usepackage{amssymb}
\usepackage{newtxtext} \usepackage{newtxmath}
\usepackage{amsmath, amssymb}
\newtheorem{problem}{Problem}
\newtheorem{example}{Example}
\newtheorem{lemma}{Lemma}
\newtheorem{theorem}{Theorem}
\newtheorem{problem}{Problem}
\newtheorem{example}{Example} \newtheorem{definition}{Definition}
\newtheorem{lemma}{Lemma}
\newtheorem{theorem}{Theorem}
\DeclareMathAlphabet{\mathcal}{OMS}{cmsy}{m}{n}


\begin{document}

$\S$ El sentimiento de vacío no consiste en no sentir nada. Tiene un
\textit{qualia} determinado, pero cuya determinación es precisamente la no
determinación, la vacuidad. En la promiscuidad obscena de mis veintes no
encontré nada más que esto. Sentía un afecto tan superficial por mis amantes que
acaso era apenas una forma de la distracción. Lo que es peor, si fui querido por
ellas, donde yo presentí la amorfa intuición de lo que no es, junto con una
viciosa y primitiva forma de interés, ellas pudieron sentir amargura. Me vaciaba
a mí mismo como un frasco lleno de humo que se destapa de repente mientras todo
a su alrededor, aunque se llena, se llena de tinieblas. Y en todo era movido por
la sed idiota del que busca dar caza a una entidad desconocida, una que, cuando
finalmente es alcanzada, aunque \textit{en cierto modo} muere, crea una copia de
sí misma en un lugar desconocido. Y todo debe comenzar de nuevo.

$\S$ Hace poco soñé que me encontraba en una mesa rodeado de gente que aprecio y
que me aprecia. Una colega científica decía «en verdad es un joven inteligente»,
la directora de mi laboratorio celebraba mis logros, y mi pareja me contemplaba
con amor diciéndome lo noble que era. Entonces presentí mi absoluta pequeñez y,
poniendo las manos sobre la cruz que pende de mi cuello, lloré con amargura.
Nadie comprendió mi llanto, pero les dije: «¡Si tan sólo hubiera hecho todo lo
bueno que hice en mi vida, pero sin que nadie supiera que fui yo el que lo
hizo!» Porque siempre cultivé la arrogancia y la vanidad, y me sentí mejor que
otros por tener hinchada la cabeza, sin darme cuenta que la única virtud está en
el corazón del hombre. Por eso a mis quince años ---lo recuerdo vivamente, y es
la memoria que más vergüenza me da--- me reí con sorna de una humilde persona
que escuché consultar por libros de filosofía, como pensando: «¡como que ésta va
a leer algo!» La vida me enseñó luego la perfección de esa mujer, como para
enseñarme una lección, pero yo era aún demasiado estúpido para aprender nada.
Vivía hinchado de mí mismo, y cuanto más se inflaba mi cabeza más se torcía mi
corazón. Adoraba que me elogien, y para colmo de males cultivaba precisamente
actividades que la gente gusta de elogiar ---la lectura, la ciencia, la música y
la poesía---. Verdaderamente digo, quisiera nunca más ser elogiado, y cada obra
de amor quisiera hacerla en secreto. Y si amo lo que hago, ¿qué puede darme un
elogio? El amor es su propio fin.

$\S$ Me gustaba repetir, en mi ignorancia, aquella sentencia del
\textit{Apocalipsis} que me enseñó mi padre; y, considerándome en cada posición
tomada firme y fuerte, decía «Dios vomita a los tibios». No sólo todo lo juzgaba
sino que me enorgullecía de hacerlo, considerándolo una virtud, porque se me
había enseñado que para ser justo hay que juzgar. Pero la etimología es
engañosa, porque \textit{justo} deriva de \textit{iūstus}, del que está de
acuerdo con la ley; y juzgar de \textit{iudicāre}, de dar un veredicto. Pero en
mi ley está escrito: \textit{Do not judge, or you too will be judged}. Porque no
es justo el que juzga, sino el que perdona y el que olvida. Y mucho me decía a
mí mismo, viendo que mi adicción al juicio me enfermaba, «¡cuánto quisiera no
juzgar a nadie!». Porque andaba solo por la vida, y a nadie dejaba que se
acerque, porque en todos veía sólo lo peor. Por eso tuve siempre pocos amigos,
aunque eso sí, excelentes, mucho más excelentes que yo. Y a todos ellos herí con
mi juicio, de una u otra forma, y a mí también.


$\S$ Y además, vida mía, me alejaste de aquellos que más quise, y de las pocas
personas que acepté en mi corazón. A mis tres amigos más hermosos los dejé en mi
ciudad natal; a Paulina la llevó la muerte; me enamoré de una joven llamada
Carla y, poco después, se fue a Buenos Aires; luego de Iraí, que se fue a
Santiago; a Selene y a Mutio, las dos personas que más amo en este mundo, las
alejé de mí por años, y a ambas por acción de mi arrogancia intemperante; a mis
hermanos, porque nuestra familia vivió rota y no quería comprenderlos; a mi
madre por mi propia ignorancia, y a mi padre por la del mundo. Por eso a mis
veintiún años consideré la muerte, y seriamente, y sopesé con sumo cuidado cada
factor relevante, y decidí cómo me quitaría la vida. Pero una tarde presentí una
imagen, que no vi con los sentidos, sino que fue como un latido interior: un
pastor rodeado de luz cuidaba de una oveja, y nada más. Y al sentir este latido
me inundó la paz del mundo. Inmediatamente hablé con un amigo cristiano, y me
dijo «fue Cristo». Y yo le dije: «no fue nadie, fui yo, y lo que importa no es
quién fue, sino que ahora no quiero morir». 

$\S$ Entonces me reconcilié con aquellos a quienes mi arrogancia había herido, y
organicé mi vida, y orienté mis energías a la ciencia y el estudio, y me sentí
sano y dinámico. Pero, aunque ya no estuve solo, seguía lleno de vanidad, y
cuanto más estudiaba, y cuanto más prosperaba en el trabajo científico, más me
envanecía.

$\S$ Y una amante, con la que obré mal, me confesó un día que al conocerme
presintió en mí una tristeza y una angustia inconmensurables. Y me sorprendió,
no por no reconocerlas, sino por pensar que las ocultaba perfectamente. Y tuvo
razón, y porque me vio desnudo de este modo no quise que me viera nunca más.
Porque quien me ve es quien puede herirme, o así razonaba entonces, sabiéndolo
sólo a medias. Y me alejé sin decirle nada. Por lo que, bajo el aparente
dinamismo de mi reencarnación, seguía siendo oscuro y escurridizo.

$\S$ Y a Paulina, a la que mencioné, ¿no le hice tanto daño? Porque a los
diecinueve años me confesó su enfermedad, y yo, siendo demasiado inmaduro e
ingenuo, no supe dimensionar su gravedad. Aquella noche ---ella estaba de visita
en Córdoba, y conversábamos bajo la sola luz de mi lámpara de estudiante--- sólo
quise darle una muestra de mi amor que perdurase en el tiempo. Y, lo recuerdo
bien, improvisé un poema en un papel; y el poema, de mala calidad, pero traido
desde el fondo de mí, decía:

\begin{verse}
Si estás caminando sola,
sola y sin nadie,
o en la luz de tu silencio
te me apagases,
una sóla cosa es cierta:
yo iré a abrazarte.


Yo iré a abrazarte,
pues soy como el aire.

Aunque una lluvia violenta
hunda la calle,
abrí todas las ventanas.
Sobre tu carne
irá mi cálido aliento
cuando lo exhale.

Irá mi aliento,
pues soy como el aire.

En el rumor de la noche,
bajo la tarde 
de grillos y de silencios,
deja que pase 
el viento y no estarás sola:
yo iré a abrazarte.

(Tanto yo sé, que soy viento,
y hay tanto que tú no sabes...)
\end{verse}

Amiga del alma, ¿te hice daño? Daría mi vida veinte veces porque vivieras otra
vez, porque pudieras ver el pulso crepitante de la luna, o abrieras los ojos
para pensar: «he despertado». Porque te amé como pude, y como quise, pero no
como necesitabas. Y porque, tras tu última internación, Selene volvió a mi vida
y la amé todavía más, y me comprometí con ella, incluso cuando unos meses antes
vos y yo nos escribíamos cartas de amor. ¿Qué sentiste al verme amándola de este
modo, a ella, que era nuestra amiga? Te quitaste la vida, y para martirio de la
mía, el último audio tuyo que poseo es el de tu voz llorando cuánto me
extrañabas, y pidiéndome que pasara tiempo a tu lado, y sollozando que ya nunca
me veías solo. Y tengo que aferrarme a duras penas a la luz de mi Cristo cuando,
por la noche, un anhelo acongojado me pide que lo escuche, sabiendo que rompería
mi corazón, si tan solo sea por volver a escuchar tu voz. El reino de los cielos
está dentro de nosotros, eso lo sé, y por eso sé que te llevo conmigo, y que
cada obra de amor que haga es tu homenaje; porque no podría vivir si vivir fuera
olvidarte.

$\S$


























\end{document}



